\documentclass[a4paper,12pt]{article}
\usepackage[ngerman]{babel}
\usepackage[utf8]{inputenc}
\usepackage[T1]{fontenc}
\usepackage{amsmath,amssymb}
\usepackage{tikz}
\usetikzlibrary{automata,positioning,arrows.meta}

\tikzset{
    ->,>=Stealth,
    node distance=2.5cm,
    every state/.style={thick,minimum size=1cm},
    initial text={},
    derivation tree/.style={
        grow=down,
        level distance=1.15cm,
        sibling distance=1.35cm,
        every node/.style={inner sep=1.5pt},
        edge from parent/.style={draw,-}
    },
}

\newcommand{\methodbox}[1]{%
\medskip
\noindent\fbox{%
    \begin{minipage}{0.94\linewidth}
    \textbf{Methodik / Definition.}\par\smallskip
    \small
    #1
    \end{minipage}%
}
\medskip\par
}

\newcommand{\examplebox}[1]{%
\medskip
\noindent\fbox{%
    \begin{minipage}{0.94\linewidth}
    \textbf{Beispiele.}\par\smallskip
    \small
    #1
    \end{minipage}%
}
\medskip\par
}

\newcommand{\taskbox}[1]{%
\medskip
\noindent\fbox{%
    \begin{minipage}{0.94\linewidth}
    \textbf{Aufgabenstellung.}\par\smallskip
    \small
    #1
    \end{minipage}%
}
\medskip\par
}

\newcommand{\bits}{\{0,1\}}

\begin{document}
\raggedright
\setlength{\parindent}{0pt}

{\LARGE\bfseries ETI -- Aufgabenblatt 4 \\ Hausaufgaben\par}
\vspace{2em}

\section*{Aufgabe 4.5}

\taskbox{
Geben Sie fuer jede der drei Ableitungsregeln fuer $S$ ein Wort der Laenge fuenf an, das von der
Grammatik generiert wird, und zeichnen Sie fuer jedes dieser Woerter einen Ableitungsbaum.
Ist die Grammatik eindeutig? Begruenden Sie Ihre Antwort.
}

\methodbox{
Ein Ableitungsbaum hat innen Nichtterminale und an den Blaettern Terminale. Die Blaetter von links
nach rechts ergeben das erzeugte Wort. Um zu zeigen, dass eine Grammatik nicht eindeutig ist,
genuegt ein einziges Wort mit zwei verschiedenen Ableitungsbaeumen.
}

\textbf{Regel $S\to 0F1$.}

Wir waehlen
\[
    00\#11.
\]
Eine Ableitung ist
\[
    S\Rightarrow 0F1\Rightarrow 00\#11.
\]

\begin{center}
\begin{tikzpicture}[
    derivation tree,
    level 1/.style={sibling distance=2.2cm},
    level 2/.style={sibling distance=0.9cm}
]
    \node {$S$}
        child {node {$0$}}
        child {node {$F$}
            child {node {$0$}}
            child {node {$\#$}}
            child {node {$1$}}
        }
        child {node {$1$}};
\end{tikzpicture}
\end{center}

\textbf{Regel $S\to N1Y$.}

Wir waehlen
\[
    0\#101.
\]
Eine Ableitung ist
\[
    S\Rightarrow N1Y\Rightarrow 0\#1Y\Rightarrow 0\#10Y\Rightarrow 0\#101.
\]

\begin{center}
\begin{tikzpicture}[
    derivation tree,
    level 1/.style={sibling distance=2.5cm},
    level 2/.style={sibling distance=1.0cm},
    level 3/.style={sibling distance=1.0cm},
    level 4/.style={sibling distance=0.8cm}
]
    \node {$S$}
        child {node {$N$}
            child {node {$0$}}
            child {node {$Y$}
                child {node {$\varepsilon$}}
            }
            child {node {$\#$}}
        }
        child {node {$1$}}
        child {node {$Y$}
            child {node {$X$}
                child {node {$0$}}
            }
            child {node {$Y$}
                child {node {$X$}
                    child {node {$1$}}
                }
                child {node {$Y$}
                    child {node {$\varepsilon$}}
                }
            }
        };
\end{tikzpicture}
\end{center}

\textbf{Regel $S\to E0Y$.}

Wir waehlen
\[
    1\#001.
\]
Eine Ableitung ist
\[
    S\Rightarrow E0Y\Rightarrow 1\#0Y\Rightarrow 1\#00Y\Rightarrow 1\#001.
\]

\begin{center}
\begin{tikzpicture}[
    derivation tree,
    level 1/.style={sibling distance=2.5cm},
    level 2/.style={sibling distance=1.0cm},
    level 3/.style={sibling distance=1.0cm},
    level 4/.style={sibling distance=0.8cm}
]
    \node {$S$}
        child {node {$E$}
            child {node {$1$}}
            child {node {$Y$}
                child {node {$\varepsilon$}}
            }
            child {node {$\#$}}
        }
        child {node {$0$}}
        child {node {$Y$}
            child {node {$X$}
                child {node {$0$}}
            }
            child {node {$Y$}
                child {node {$X$}
                    child {node {$1$}}
                }
                child {node {$Y$}
                    child {node {$\varepsilon$}}
                }
            }
        };
\end{tikzpicture}
\end{center}

\textbf{Eindeutigkeit.}
Die Grammatik ist nicht eindeutig. Das Wort $00\#11$ hat mindestens zwei verschiedene Ableitungen:
\[
    S\Rightarrow 0F1\Rightarrow 00\#11
\]
und
\[
    S\Rightarrow N1Y\Rightarrow 0Y\#1Y
      \Rightarrow 00\#1Y
      \Rightarrow 00\#11.
\]
Da schon die erste verwendete Regel fuer $S$ verschieden ist, sind die zugehoerigen
Ableitungsbaeume verschieden. Folglich ist $G$ nicht eindeutig.

\section*{Aufgabe 4.6}

\taskbox{
Geben Sie fuer jedes Nonterminal $v\in V$ der Grammatik $G$ eine moeglichst einfache Beschreibung
der von $v$ generierten Sprache $L(G_v)$ an.
}

\methodbox{
Fuer ein Wort $w\in\bits^*$ bezeichnet $w^R$ die Umkehrung von $w$.
Mit $\overline{w}$ bezeichnen wir das bitweise Komplement, also den Tausch von $0$ und $1$.
}

Die von den Nichtterminalen erzeugten Sprachen sind:
\[
\begin{aligned}
L(X) &= \bits,\\
L(Y) &= \bits^*,\\
L(F) &= \{\,w\#(\overline{w})^R \mid w\in\bits^*0\,\},\\
L(N) &= \{\,u0v\#w \mid u,v,w\in\bits^*,\ |u|=|w|\,\},\\
L(E) &= \{\,u1v\#w \mid u,v,w\in\bits^*,\ |u|=|w|\,\},\\
L(S) &= \{\,u\#v \mid u,v\in\bits^*,\
            \exists i\in\{1,\ldots,\min(|u|,|v|)\}: u_i\ne v_i\,\}.
\end{aligned}
\]
In der letzten Zeile sind die Woerter $u$ und $v$ von links nach rechts nummeriert, also
$u=u_1\cdots u_{|u|}$ und $v=v_1\cdots v_{|v|}$.

\medskip
\textbf{Begruendung.}
$X$ erzeugt genau ein Bit. $Y$ erzeugt durch wiederholtes Voranstellen eines Bits und anschliessendes
Beenden mit $\varepsilon$ genau alle Binaerwoerter.

\medskip
Bei $F$ ist die Basis $0\#1$. Jede rekursive Regel fuegt links ein Bit und rechts das komplementaere
Bit an. Weil rechts aussen angefuegt wird, steht rechts die umgekehrte komplementierte linke Seite.
Also gilt
\[
    L(F)=\{\,w\#(\overline{w})^R \mid w\in\bits^*0\,\}.
\]

\medskip
Bei $N$ ist die Basis $0Y\#$. Jede Anwendung von $XNX$ fuegt links und rechts je ein beliebiges Bit
hinzu. Deshalb entstehen genau die Woerter $u0v\#w$ mit gleich langen Randstuecken $u$ und $w$.
Analog entsteht bei $E$ dieselbe Form, nur mit einer ausgezeichneten $1$ statt der ausgezeichneten $0$.

\medskip
Fuer $S$ markieren die Regeln $N1Y$ und $E0Y$ eine Position, an der die linke und rechte Seite des
$\#$ verschieden sind. Genauer erzeugt $N1Y$ eine Position mit $0$ links und $1$ rechts; $E0Y$
erzeugt eine Position mit $1$ links und $0$ rechts. Die Regel $0F1$ erzeugt nur Woerter, die bereits
in diese Beschreibung fallen, denn links beginnt das Wort mit $0$ und rechts beginnt es mit $1$.

\section*{Aufgabe 4.7}

\taskbox{
Es sei $L=\{a^ib^j\mid i,j\in\mathbb{Z}_{\ge 0},\ i\ne j\}$. Beweisen Sie, dass $L$ nicht regulaer ist.
}

\methodbox{
Regulaere Sprachen sind abgeschlossen unter Differenz. Wenn $A$ und $B$ regulaer sind, dann ist
auch $A\setminus B=A\cap \overline{B}$ regulaer.
}

Wir beweisen die Aussage per Widerspruch. Angenommen, $L$ waere regulaer.
Die Sprache
\[
    a^*b^*
\]
ist regulaer. Wegen der Abschlusseigenschaften regulaerer Sprachen waere dann auch
\[
    a^*b^*\setminus L
\]
regulaer.

Nun gilt aber
\[
    a^*b^*\setminus L
    =
    \{a^ib^j\mid i,j\in\mathbb{Z}_{\ge 0},\ i=j\}
    =
    \{a^nb^n\mid n\in\mathbb{Z}_{\ge 0}\}.
\]
Diese Sprache ist nicht regulaer. Das ist ein Widerspruch. Also kann $L$ nicht regulaer sein.

\section*{Aufgabe 4.8}

\taskbox{
Geben Sie eine kontextfreie Grammatik an, die
$L=\{a^ib^j\mid i,j\in\mathbb{Z}_{\ge 0},\ i\ne j\}$ generiert. Begruenden Sie zudem die Korrektheit.
}

\methodbox{
Wir zerlegen $L$ in die beiden Faelle $i>j$ und $i<j$. Danach erzeugt eine Startregel die Vereinigung
dieser beiden kontextfreien Sprachen.
}

Eine passende Grammatik ist
\[
G_L=(V,\{a,b\},R,S)
\]
mit
\[
V=\{S,A,B,C,D\}
\]
und Regeln
\[
\begin{aligned}
S &\to A \mid B,\\
A &\to aAb \mid C,\\
C &\to aC \mid a,\\
B &\to aBb \mid D,\\
D &\to bD \mid b.
\end{aligned}
\]

\textbf{Korrektheit.}
Das Nichtterminal $A$ erzeugt zuerst gleich viele $a$- und $b$-Paare durch Regeln der Form
$aAb$ und beendet dann mit mindestens einem zusaetzlichen $a$ aus $C$. Also gilt
\[
    L(A)=\{a^ib^j\mid i>j\}.
\]
Insbesondere werden auch Woerter der Form $a^i$ mit $i>0$ erzeugt, indem keine $b$-Paare benutzt
werden.

\medskip
Analog erzeugt $B$ zuerst gleich viele $a$- und $b$-Paare und beendet dann mit mindestens einem
zusaetzlichen $b$ aus $D$. Also gilt
\[
    L(B)=\{a^ib^j\mid i<j\}.
\]
Insbesondere werden auch Woerter der Form $b^j$ mit $j>0$ erzeugt.

\medskip
Die Startregel $S\to A\mid B$ bildet die Vereinigung beider Faelle:
\[
    L(G_L)=L(A)\cup L(B)
    =
    \{a^ib^j\mid i,j\in\mathbb{Z}_{\ge 0},\ i\ne j\}.
\]

\end{document}
